\setlength{\absparsep}{18pt}

\begin{resumo}[RESUMO]
Este trabalho propõe a arquitetura e o plano de implementação de um hiperinstrumento musical gestual denominado Aerochord, que utiliza visão computacional por câmera RGB comum para gerar eventos MIDI 2.0 com resolução estendida, buscando democratizar a expressão musical para usuários sem formação técnica em instrumentos e pessoas com mobilidade reduzida. Inicialmente, foi realizada a revisão bibliográfica em interação humano-computador e computação musical, seguida de análise crítica do estado da arte em interfaces gestuais, o que evidencia a lacuna entre acessibilidade, latência e expressividade profissional. Com base nessa investigação, foram definidos os requisitos funcionais e não funcionais para o Aerochord, incluindo processamento de vídeo em tempo real, detecção de pontos de referência das mãos via técnicas de visão computacional, mapeamentos gesto–som preliminares e comunicação via MIDI 2.0 com plano de contingência para modo legado. Foi proposta uma arquitetura modular em C++ em conjunto com o \textit{framework} JUCE, organizada em pipelines buscando minimizar o \textit{jitter}, uso de filas sem bloqueio, \textit{handshake} por \textit{MIDI Capability Inquiry} e inserção de \textit{timestamps} em pacotes universais, de modo a garantir resposta sonora coerente a gestos do performer. Planeja-se, na fase seguinte, avaliação experimental com medições objetivas de latência e precisão de reconhecimento, além de métodos subjetivos de usabilidade e carga cognitiva, incorporando cuidados de privacidade e adaptabilidade a diferentes perfis de usuário e cenários de uso. Esta etapa conceitual estabelece fundamentos robustos para a implementação do protótipo, oferecendo diretrizes técnicas e metodológicas para prototipagem em tempo real e testes experimentais. Espera-se que o Aerochord contribua para o campo de interação humano-computador e computação musical ao apresentar um modelo de hiperinstrumento acessível, responsivo e expressivo, servindo de base para pesquisas futuras em mapeamentos adaptativos e ambientes imersivos.

\textbf{Palavras-chave}: interfaces gestuais, computação musical, MIDI 2.0, democratização musical, visão computacional, hiperinstrumentos, acessibilidade, baixa latência.
\end{resumo}